\section{From minimizers to finite switching patterns}\label{sec:outline-periods}

An admissible control may take any value in the control triangle. A global minimizer, however, can
be represented by finitely many pieces of constant vertex control. We establish this
representation and classify its switching patterns. Excluding the longer patterns will leave at
most five pieces.

\subsection{Finiteness and normality}\label{sec:outline-reduction}

The \emph{Pontryagin maximum principle} gives necessary conditions for an optimal control. It
introduces auxiliary matrix-valued functions, called \emph{costates}, that measure the response of
the cost to variations of the trajectory. Together with the state and control, they form a
\emph{Pontryagin extremal} (\cref{def:pmp-extremal}). A scalar multiplier records the contribution
of the cost. The extremal is \emph{normal} when this multiplier is nonzero; it can then be
normalized to $-1$.

These necessary conditions do not by themselves give a finite switching description. At certain
states and costates, called the \emph{singular locus}, they fail to distinguish the available
controls. Such behaviour can lead to infinitely many switches in finite time. Our first estimate
uses the bound $\cost\le\Aoct<\pi$ to keep a minimizing extremal away from this locus
(\cref{prop:clearance}).

Away from the singular locus, one of the three vertex controls is locally excluded. The problem
then reduces locally to an edge of the control triangle. An edge minimizer has at most four
pieces (\cref{prop:edge-bound}). Covering the time interval by finitely many such neighbourhoods
gives a finite number of switches. After joining equal neighbouring pieces and choosing a switch
as the starting time, the minimizer has a \emph{vertex representation}: a finite sequence of
constant vertex controls with unequal neighbouring labels, including across the closing seam.
Finally, the closing conditions rule out a zero cost multiplier.

The resulting object is a \emph{normal vertex period}: an admissible trajectory with a vertex
representation and a normal Pontryagin extremal. It is \emph{competitive} when its cost is at
most $\Aoct$. The finite normal representation theorem (\cref{thm:finite-normal}) gives such a
period for every global minimizer. It has the same cost as the original trajectory and stays in
$\hstarstar$; only its starting time has changed.

The equality case requires this representation of the original disk's motion, including its change
of starting time. The precise reduction is proved in \cref{sec:reduction}.

\subsection{Data of a single piece}\label{sec:outline-pieces}

Let $N$ be the number of pieces. We write $\tau_k$ for their switching times and
$i_k\in\mathbb Z/3\mathbb Z$ for their labels, where labels are read cyclically as $0,1,2$.
The closing condition shifts every label by one after a period, so $i_{k+N}=i_k+1$.
For the octagon, the four labels are $2,0,1,2$; the next label is $0$, obtained by adding one to
the first label.

The state has three positive coordinates $p_0(X),p_1(X),p_2(X)$, called \emph{margins}, one for
each vertex control (\cref{lem:margins}). On a piece with label $i$, the coordinate $p_i(X)$ is
constant. It is the \emph{active margin} of that piece. Write $a_k$ for this constant on piece
$k$. It measures the scale of the state's motion during that piece. The other two margins, the
frame and the costate all have explicit evolution formulas.

A piece can therefore be specified by its starting data and one positive duration. The state and
frame must close after all the pieces have been composed. In these variables, closure is a finite
system of equations, and the maximum principle becomes a finite system of inequalities.
\Cref{sec:periods} derives the formulas and their consequences for the case arguments.

\subsection{The five switching patterns}\label{sec:outline-signs}

A switch is \emph{forward} if $i_{k+1}=i_k+1$ and \emph{reverse} if
$i_{k+1}=i_k-1$. Denote its sign by $+$ or $-$, respectively. These are the only possibilities
because consecutive labels are unequal. A \emph{run} is a maximal consecutive block of switches
with the same sign. The signs repeat after a period, so runs are understood cyclically. A period
is \emph{mixed} if it has both signs.

\begin{definition}[The five classes]\label{def:classes}
A vertex period is in class
\begin{itemize}
\item $\cls{R}$ if every switch is reverse;
\item $\cls{F}$ if every switch is forward;
\item $\cls{M3}$ if it is mixed and some reverse run has length other than two;
\item $\cls{M1}$ if it is mixed, every reverse run has length two and every forward run has
  length one, so its cyclic sign sequence is $(--+)^m$ for some positive integer $m$;
\item $\cls{M2}$ if it is mixed, every reverse run has length two and some forward run has
  length at least two.
\end{itemize}
\end{definition}

\begin{proposition}[Partition]\label{prop:partition}
Every vertex representation is in one of the classes $\cls{R}$, $\cls{F}$, $\cls{M3}$,
$\cls{M1}$, $\cls{M2}$.
\end{proposition}

In a mixed period, the maximum principle excludes reverse runs whose length differs from two. The
remaining patterns require comparisons with the cost of a complete period.

\subsection{A conserved quantity and the cost budget}\label{sec:outline-casimir}

The costates supply a quantity that links these local and global conditions. For the first
costate matrix $\Lambda_1(t)$, conjugation along the trajectory preserves its determinant. We
use the scaled invariant
\[
  \gamma^2=\frac{\det\Lambda_1(0)}{12},
\]
called the \emph{Casimir}. This notation does not assume positivity. Every mixed normal period has $\gamma^2>0$ (\cref{prop:ellipticity}).

The frame must turn through $\pi/3$ in a period. The cost of this turning gives an upper bound
for the Casimir: every competitive mixed period satisfies
\[
  \gamma^2<\frac5{27}.
\]
The cost--Casimir inequality and this consequence are proved in
\cref{prop:cost-casimir,cor:casimir-cap}. The entirely forward case uses a sharper bound of the
same kind.

The invariant can also be computed from the data adjacent to a single piece. In two of the case
arguments, these local data force a value of $\gamma^2$ exceeding the bound imposed by the cost. A
third case compares the turning accumulated at successive pieces with that required by closure.

\subsection{The bound on the number of pieces}\label{sec:outline-five}

\begin{theorem}[At most five pieces]\label{thm:at-most-five}
A competitive normal vertex period whose state stays in the compact domain $\hstarstar$ has at
most five pieces.
\end{theorem}

\Cref{sec:outline-classes} proves this theorem by excluding the five possible classes of longer
periods. The remaining periods admit a direct comparison with the octagon.